% [原文 p.127]
在构造作用量时，我们需要保证 $P$ 的运动方程给出 $P=d_A^{\star}\varphi$。一个合适的作用量是 BRST-正合的：
\[ S = -\frac{1}{2e^2}\delta_Q\int_\Sigma\Tr\left(\ii\zeta\wedge\star\bar\cF + \ii\beta\wedge\star(P-2d_A^{\star}\varphi)\right) \ . \]

% [Tutor Note: p.127 补出省略论证：为什么 $P$ 的运动方程是 $P=d_A^{\star}\varphi$。作用量的 $P$ 依赖只来自 $\delta_Q(\ii\beta\wedge\star(P-2d_A^{\star}\varphi))$ 一项；用 $\delta_Q\beta=\ii P$、$\delta_QP=0$ 得该项为 $-\ii P\wedge\star(P-2d_A^{\star}\varphi)$，再乘上外层的 $-\frac{1}{2e^2}$ 与 $\Tr$，对 $P$ 变分即得 $P=d_A^{\star}\varphi$。这也说明引入 $P$ 的代价：作用量中出现了 $\delta_Q\beta=\ii P$ 带来的非正合部分，而 $P$ 在壳上被消去。]

耦合常数 $e^2$ 只作为 BRST-正合项的系数出现，因此拓扑关联函数不依赖于它，半经典近似是精确的。理论的拓扑性也是显然的，因为度规只通过 BRST-正合项进入。

\para{B-型 2d 规范理论的可观测量}

通常，局部可观测量定义为场的 BRST 不变且规范不变的标量函数，模去 BRST 变换。在目前的情形下，没有这种非平凡的局部可观测量。然而，存在非平凡的、BRST 不变的局部无序算子，它们是通过允许场中出现某些奇异性来定义的。例如，可以要求联络 $\cA$ 在插入点周围具有非平凡的和乐。这样的局部算子类似于四维规范理论中的 Gukov-Witten 面算子。更系统地说，要确定允许哪些种类的局部算子，可以把 2d 规范理论在圆周上约化，并研究所得 1d TFT 的态空间。在目前的情形下，这个 1d TFT 是一个靶空间为 $G_\C$ 的带规范场的 $\sigma$-模型。从 2d 的观点看，靶空间参数化 $\cA$ 的和乐。BRST 不变的波函数是 $G_\C$ 上在共轭下不变的全纯函数，即 $G_\C$ 的特征标。更一般地，可以考虑不可归一化的波函数，例如支撑在 $G_\C$ 的闭 $G_\C$-不变复子流形上的 $\delta$-函数。举例来说，恒等算子可以看成一个支撑在单位元处的 $\delta$-函数，而 Gukov-Witten 型的局部算子是支撑在 $G_\C$ 中闭共轭类上的 $\delta$-函数。

% [Tutor Note: p.127 术语首现与原文笔误。术语：disorder operator 译“无序算子”（与“有序算子”相对，后者由场的正则表达式给出；无序算子的定义方式是允许场在插入点附近有指定奇异性）。原文此处还有一处衍字：“one can reduce the 2d gauge theory theory on a circle”（theory 重复），正确形式为 “the 2d gauge theory on a circle”。这段“约化到圆周上”的方法是确定局部算子谱的一般手段，与第 5 章把 4d 理论在 $S^1$ 上紧化的做法同源。]

还有 BRST 不变且规范不变的线可观测量，其中最基本的是复 BRST 不变联络 $\cA$ 的 Wilson 线算子。为定义它们，需要选取 $G$ 的一个有限维分次表示 $V$，并考虑 $\cA$ 在表示 $V$ 中的和乐。

\para{B-型 2d 规范理论的膜范畴}

这个 2d TFT 的膜范畴是 $G$ 的有限维分次表示的范畴。为看出这一点，考虑规范场的 Neumann 边界条件，即让 $\cA$ 在边界上的限制保持自由，并要求 $\star\varphi$ 的限制为零。BRST 不变性于是要求 $\zeta$ 与 $\star\lambda$ 的限制在边界上为零。由于边界上的规范场 $\cA$ 不受约束且 BRST 不变，我们可以把 $G$ 的任意有限维分次表示 $V$ 耦合到它上面。也就是说，我们可以在路径积分中包含 $\cA$
% [原文 p.128]
在表示 $V$ 中的和乐。因此边界条件自然地由 $G$ 的表示来标记。给定任意两个不可约表示 $V_1$ 与 $V_2$，只有当 $V_1$ 与 $V_2$ 同构时才能在它们之间构造交接点（因为在 Neumann 边界上没有非平凡的 BRST 不变局部算子）。进一步，若 $V_1\simeq V_2$，则它们之间的态射空间是 $\Hom_G(V_1,V_2)$（出于同样的理由）。

% [Tutor Note: p.128 补出省略论证。原文说“出于同样的理由”给出了两个结论，其共同依据是：Neumann 边界上只有平凡的 BRST 不变局部算子，因此从 $V_1$ 到 $V_2$ 的交接点只能由 $G$-等变映射给出——若 $V_1\not\simeq V_2$ 则不存在这样的映射（不可约表示之间的 $G$-等变映射只能是零），而 $V_1\simeq V_2$ 时其空间恰为 $\Hom_G(V_1,V_2)$。这与“边界条件由表示标记”一起，说明膜范畴就是 $G$ 的有限维分次表示范畴。]

\section{带规范场的 2d B-模型}

本节描述如何把一个 B-型 2d 规范理论耦合到一个 B-模型上。我们将证明，所得 2d TFT 的膜范畴与凝聚层的等变导出范畴密切相关。

设 $X$ 是一个 Calabi-Yau 流形（即带全纯体积形式的 Kähler 流形），它容许一个同时保持 Kähler 结构与全纯体积形式的 $G$-作用。B-模型物质扇区的玻色场是丛
\[ E = P\times_G X \]
的一个截面 $\Phi$，该丛由全纯 $G$-作用从规范丛 $P$ 相伴得到。B-模型物质扇区的费米子场由如下取值于垂直子丛的形式组成：
\begin{gather*}
\eta \in \Omega^0(\Sigma, \Phi^*V^{0,1}E) \\
\theta \in \Omega^0(\Sigma, \Phi^*(V^{1,0}E)^\vee) \\
\rho \in \Omega^1(\Sigma, \Phi^*V^{1,0}E) \ .
\end{gather*}

% [Tutor Note: p.128 术语首现。$VE\subset TE$ 记 $E$ 的垂直切丛，$V^{1,0}E$、$V^{0,1}E$ 是其按 $X$ 的复结构分解出的 $(1,0)$ 与 $(0,1)$ 部分；$\theta$ 取值于 $(V^{1,0}E)^\vee$（对偶）而不是 $V^{1,0}E$ 本身，这一点在写作用量 $\theta_i\rho^j\ldots$ 的指标配对时很关键：$\theta_i$ 是“下指标”、$\rho^i$ 是“上指标”，二者的收缩给出标量。若把 $\theta$ 误写成同一切丛的形式，指标的升降与 $R^i{}_{jk\bar l}$ 的缩并都不再自洽。]

BRST 变换由通常 B-模型变换的规范协变化版本给出：
\begin{align*}
\delta_Q\phi^i &= 0, \\
\delta_Q\phi^{\bar i} &= \eta^{\bar i}, \\
\delta_Q\eta^{\bar i} &= 0, \\
\delta_Q\theta_i &= 0, \\
\delta_Q\rho^i &= \cD\phi^i
\end{align*}
其中 $\cD\phi^i = d\phi^i + V^i(\cA)$ 是关于复化联络 $\cA$ 的规范协变导数，而 $V$ 是对应于 $G$-作用的全纯向量场（取值于 $\fg^*$）。由于 $\cA$
% [原文 p.129]
是 BRST 不变的，这些 BRST 变换仍然满足 $\delta_Q^2=0$。

% [Tutor Note: p.128–129 可验证点。“由于 $\cA$ 是 BRST 不变，故 $\delta_Q^2=0$ 仍成立”可以直接核对：$\delta_Q^2\phi^{\bar i}=\delta_Q\eta^{\bar i}=0$，而 $\delta_Q^2\rho^i=\delta_Q(\cD\phi^i)=\cD(\eta^i)+\delta_Q(V^i(\cA))$，后者为零正是因为 $\delta_Q\cA=0$（见 p.126 的批注）与 $V^i$ 全纯。另注：$\cD$ 与 $V^i(\cA)$ 中的 $\cA$ 都是“协变化”的产物——普通 B-模型的 $\delta_Q\rho=d\phi$ 在这里被替换成 $\cD\phi$，与非规范情形相比多出 $V^i(\cA)$ 项。]

为构造 BRST 不变的作用量，我们从 B-模型的通常作用量出发并把所有导数规范协变化。协变化后的作用量不是 BRST 不变的，但这可以通过加入一个正比于 $\theta_iV^i(\zeta)$ 的新项来修正，其中 $\zeta$ 是 B-型 2d 规范理论的费米子、$ad\,P$-值 2-形式。完整的物质作用量为
\[ S = \int_\Sigma\delta_Q\left(g_{i\bar j}\rho^i\wedge\star\bar\cD\phi^{\bar j}\right) + \int_\Sigma\left(-\ii\theta_iV^i(\zeta) + \theta_i\cD\rho^i + \frac{1}{2}R^i{}_{jk\bar l}\theta_i\rho^j\rho^k\eta^{\bar l}\right) \ . \]
这里 $g$ 是 Kähler 度规，$R$ 是它的曲率张量，而 $\rho$ 的协变导数同时包含 Levi-Civita 联络与规范联络：
\[ \cD\rho^i = d\rho^i + \Gamma^i_{jk}d\phi^j\rho^k + \nabla_jV^i(\cA)\rho^j, \qquad \nabla_jV^i = \partial_jV^i + \Gamma^i_{jk}V^k \ . \]
协变导数 $\bar\cD\phi^{\bar j}$ 的定义使得作用量的玻色部分正定：
\[ \bar\cD\phi^{\bar j} = d\phi^{\bar j} - V^{\bar j}(\bar\cA), \qquad \bar\cA = A - \ii\varphi = -\cA^\dagger \ . \]

% [Tutor Note: p.129 可验证的符号。$\bar\cA=A-\ii\varphi=-\cA^\dagger$ 的三处写法互相印证：$\cA=A+\ii\varphi$，而 $A$ 反厄米（$A^\dagger=-A$）、$\varphi$ 作为 $ad\,P$-值 1-形式也反厄米（$\varphi^\dagger=-\varphi$），故 $\cA^\dagger=A^\dagger-\ii\varphi^\dagger=-A+\ii\varphi$，于是 $-\cA^\dagger=A-\ii\varphi=\bar\cA$。核对时注意 $(\ii\varphi)^\dagger=-\ii\varphi^\dagger$ 中的两个符号，否则会得到 $\bar\cA=\cA$ 的错误结论。另：$\bar\cD$ 中 $V^{\bar j}(\bar\cA)$ 取负号，与 $\cD$ 中 $+V^i(\cA)$ 相反，这是为了让玻色部分正定（两者互为共轭）。]

\para{带规范场的 2d B-模型的膜范畴}

既然靶空间为 $X$ 的 B-模型的膜范畴是 $D^b(\operatorname{Coh}(X))$，带规范场的 B-模型的膜范畴的一个自然猜想是 $D^b_{G_\C}(\operatorname{Coh}(X))$。下面我们将描述对应于 $X$ 上等变全纯向量丛复形的边界作用量的一个构造。设 $E$ 是 $X$ 上的分次复向量丛，具有全纯结构 $\bar\partial_E:E\to E\otimes\Omega^{0,\bullet}(X)$、$(\bar\partial_E)^2=0$，以及一个全纯的 1 次自同态 $T:E\to E$，满足 $\bar\partial_ET=0$、$T^2=0$。为写出具体的边界作用量，我们将假定还给定了 $E$ 的每个分次分支上的一个厄米度规，使得 $\bar\partial_E$ 给出 $E$ 上的一个联络 $\nabla^E$。我们把相应的联络 1-形式记作 $\omega$，其曲率记作 $F^E$。我们假定给定 $X$ 上 $G$-作用到 $E$ 全空间上 $G$-作用的一个提升，它纤维线性，并与 $\bar\partial_E$、$T$ 以及厄米度规相容。无穷小地，李代数 $\fg$ 如下作用在 $E$ 的截面 $s$ 上：
\[ (f,s)\mapsto f(s) = V^i(f)\nabla^E_i s + V^{\bar i}(f)\nabla^E_{\bar i}s + R(f)s, \qquad f\in\fg \ . \]
这里 $\nabla^E=d+\omega$，而 $R$ 是一个 0 次丛态射 $R:E\to E\otimes\fg^*$。$G$-作用与 $\nabla^E$ 交换这一条件蕴含
\[ \nabla^ER = \iota_VF^E \ . \]

% [Tutor Note: p.129 术语。$D^b_{G_\C}(\operatorname{Coh}(X))$ 读作“$X$ 上 $G_\C$-等变凝聚层的导出范畴”；下标 $G_\C$ 表示等变结构取复化群（与 §A.2 中 $\cA$ 的复化一致）。$\Omega^{0,\bullet}(X)$ 中的 $\bullet$ 是 Dolbeault 复形次数的占位符（不要写成点号 $\cdot$），所以 $\End(E)\otimes\Omega^{0,\bullet}(X)$ 是一个分次对象、其元素可按总次数分解。$\iota_V$ 是内乘（interior product），$\iota_VF^E$ 取值为 $\End(E)$ 中元素。]

% [原文 p.130]
$G$-作用与 $T$ 交换这一条件蕴含
\[ V^i\nabla^E_iT + [R,T] = 0 \ . \]
现在考虑拉回丛 $\Phi^*E$ 上如下依赖于场的联络 1-形式：
\[ \cN = \omega_id\phi^i + \omega_{\bar i}d\phi^{\bar i} - R(\cA) + \rho^i\eta^{\bar j}F^E_{i\bar j} + \rho^i\nabla^E_iT \ . \]
经过一番计算可以验证，其 BRST 变分满足
\[ \delta_Q\cN = d(\omega_{\bar i}\eta^{\bar i} + T) + [\cN, \omega_{\bar i}\eta^{\bar i} + T] \ . \]
因此其和乐的超迹是 BRST 不变的，可以作为相应路径积分中的边界权重因子。按定义，边界作用量是边界权重因子取负对数。

% [Tutor Note: p.130 补出省略论证。$\cN$ 的构造把四类信息拼在一起：$\omega_id\phi^i+\omega_{\bar i}d\phi^{\bar i}$ 是 $\nabla^E$ 沿 $\Phi$ 的拉回、$-R(\cA)$ 是等变性的修正、$\rho^i\eta^{\bar j}F^E_{i\bar j}$ 与 $\rho^i\nabla^E_iT$ 是费米子双线性项。原文说 “with some work one can check”，所验证的其实是 $\delta_Q\cN$ 具有 $d(\ldots)+[\cN,\ldots]$ 的“规范变换”形状，这正是超迹和乐 BRST 不变的原因（$\delta_Q$ 作用在联络上等于一次规范变换，而超迹在规范变换下不变）。]

让我们考虑在此类权重因子存在时的一个鬼数为零的边界可观测量 $\cO$。它是 $\End(E)$ 中依赖于场 $\Phi$、$\eta$ 且总次数为零的一个元素。更不变地说，可以把它视为 $\End(E)\otimes\Omega^{0,\bullet}(X)$ 的一个截面。在 $\cO$ 存在时，边界权重因子的 BRST 变分正比于
\[ \eta^{\bar i}\nabla^E_{\bar i}\cO + [T,\cO] \ . \]
因此 BRST 不变的边界可观测量是被 $\bar\partial_E$ 湮灭并与 $T$ 交换的 $\End(E)\otimes\Omega^{0,\bullet}(X)$ 的截面。进一步，一个 BRST 不变的 $\cO$ 是规范不变的，当且仅当它满足
\[ V^i\nabla^E_i\cO + [R,\cO] = 0 \ . \]
这些条件合起来意味着 $\cO$ 表示等变复形 $(E,\bar\partial_E,T)$ 作为 $D^b_{G_\C}(\operatorname{Coh}(X))$ 中对象的一个自同态。也很容易看出，这样的可观测量 $\cO$ 是某个规范不变可观测量的 BRST 变分，当且仅当它同伦于零。在某些情形下，这意味着我们所构造的这类带规范场的 B-模型的膜范畴等价于 $D^b_{G_\C}(\operatorname{Coh}(X))$。这发生在 $X$ 上任一 $G$-等变凝聚层都有由 $G$-等变全纯向量丛给出的 $G$-等变消解的时候。这样的 $X$ 被说成具有 $G$-消解性质。这样的 $X$ 的一个例子是带有 $G$ 的线性作用的 $\C^n$，或者更一般地，带有 $G$ 的仿射作用的光滑仿射簇。注意，对一般的复流形 $X$，即使 $G$ 是平凡的，消解性质也可能失效。但当 $G$ 平凡时补救办法是已知的：必须把全纯向量丛的复形换成 $X$ 的 Dolbeault DG-代数上更一般的 DG-模 [30]。这些更一般的 DG-模也会自然地
% [原文 p.131]
从物理的观点出现 [27,36]。我们预期，对任何在 $X$ 上有复解析作用的复李群 $G$，$X$ 上的 $G$-等变凝聚层都有由这些更一般的 DG-模给出的 $G$-等变消解。这将蕴含带规范场的 B-模型的 B-膜范畴等价于 $D^b_{G_\C}(\operatorname{Coh}(X))$。

% [Tutor Note: p.130–131 术语与补出论证。“$G$-resolution property”（$G$-消解性质）指 $X$ 上每个 $G$-等变凝聚层都有由 $G$-等变全纯向量丛构成的消解；原文只在给出这个定义后立即用它，其用途是：有消解时，向量丛复形（即 §A.3 前面构造的 $(E,\bar\partial_E,T)$）足以"跑遍"整个导出范畴，于是膜范畴的构造给出完整的等价。原文没有解释“$G$ 平凡时如何补救”，只给了文献 [30]：其思路是把“向量丛复形”换成为“在 Dolbeault DG-代数上取模”的 DG-模（Dolbeault 复形上的模结构自带足够多的投射对象），从而绕开向量丛消解不存在的问题。最后一句的“预期”是作者的猜想，不是已证定理，译文忠实保留“预期”二字。]

% [Tutor Note: p.131 未解决的问题提示。原文这里的逻辑层次要分清：已证的是“若 $X$ 具有 $G$-消解性质，则膜范畴 $\simeq D^b_{G_\C}(\operatorname{Coh}(X))$”；而“任一复李群复解析作用都有这样的 DG-模消解”只是作者的预期，故末尾的等价性在原文中也只是条件性的结论。阅读第 6 章引用本附录的结论时应记住这一条件性。]
